\par\noindent\textbf{(1)} 最小耦合Hamiltonian H=(p+eA)²/(2m*)，机械动量π=p+eA满足[πx,πy]=−iℏeB（符号由规范约定）。定义升降算符使H=ℏωc(a†a+1/2)，ωc=eB/m*>0，故En=ℏωc(n+1/2)，n=0,1,…。
\par\noindent\textbf{(2)} 磁长度ℓ\_B=√(ℏ/(eB))，每个Landau能级在面积A中允许轨道中心数Nφ=A/(2πℓ\_B²)=eBA/h，即穿过样品的磁通BA除以磁通量子h/e。边缘与有限面积导致小修正，但体态简并度与样品形状无关。
\par\noindent\textbf{(3)} 若ν个Landau能级完全填充，粒子密度n=νeB/h。Středa关系σxy=e(∂n/∂B)\_μ（电荷符号决定Hall电导取向），给Hall电导幅值|σxy|=νe²/h。亦可由ν条手性边缘通道各贡献e²/h解释。
\par\noindent\textbf{(4)} B→0时ωc→0而每级简并度∝B减小，许多能级趋连续，回到自由二维态密度；观察稳定的整数量子Hall平台还需要体态能隙与无序局域化、合适化学势，理想无无序模型中平台随磁场并非任意宽。
